\section{Régulation et algorithmes de traitement}

\subsection{Positionnement}

Le scanner n'intègre pas de boucle de régulation asservie au sens classique
sur la mécanique : il n'y a pas de capteur continu qui corrige en temps réel
la position du plateau. Il comporte néanmoins plusieurs \emph{problématiques de
contrôle} (position angulaire, exposition caméra, synchronisation
moteur/laser/caméra) et surtout une \emph{chaîne d'algorithmes de traitement}
qui conditionne directement la qualité du résultat.

Cette section détaille ces algorithmes, leur \textbf{évolution depuis le rapport
intermédiaire} et la \textbf{raison d'être de chacun}. Là où le chapitre
Software (\autoref{fig:archi-composants}) décrit l'architecture logicielle,
on s'intéresse ici au « pourquoi » et au « comment » de chaque traitement, du
pilotage des actionneurs jusqu'au maillage final.

\subsection{Contrôle du plateau rotatif (boucle ouverte)}

Le plateau est piloté en \textbf{boucle ouverte} via le driver \texttt{DM320T} :
le logiciel envoie un nombre de pas connu et reconstruit la position angulaire
par comptage logiciel. Le tour complet est échantillonné sur 100 positions
(moteur NEMA\,17, 200\,pas/tour, $1{,}8^\circ$/pas).

\paragraph{Justification du choix.}
\begin{itemize}
  \item la vitesse de rotation est lente et la charge reste modérée~;
  \item le moteur NEMA\,17 offre une marge de couple suffisante pour les objets
        visés, et aucun blocage ni perte de pas notable n'a été observé pendant
        les scans~;
  \item l'ajout d'un codeur rotatif (boucle fermée) n'est pas prioritaire tant
        que les erreurs dominantes proviennent de l'occultation, de la
        calibration et du maillage, et non du positionnement angulaire.
\end{itemize}

La limite résiduelle n'est donc pas le moteur mais le \emph{glissement} de
certaines pièces sur le plateau, traité côté mécanique (surface adhérente et
centrage).

\subsection{Régulation de l'exposition (boucle fermée)}

C'est la seule véritable \textbf{boucle fermée} du système. Au rapport
intermédiaire, l'exposition était fixée manuellement, ce qui rendait
l'extraction de la ligne laser sensible à la lumière ambiante et à la couleur
des pièces. La version finale introduit une \textbf{calibration automatique de
l'exposition} exécutée avant chaque scan, indépendamment pour chaque caméra.

\paragraph{Algorithme.}
\begin{enumerate}
  \item échantillonner quelques angles du plateau, laser allumé~;
  \item mesurer la \emph{largeur médiane} de la ligne laser détectée~;
  \item comparer à une largeur cible (de l'ordre de 6\,pixels) à une tolérance
        d'environ 1\,pixel~;
  \item corriger le temps d'exposition dans une plage bornée
        (200 à 60\,000\,$\mu$s) et recommencer jusqu'à convergence (quelques
        itérations).
\end{enumerate}

\paragraph{Raison d'être.}
La largeur de la ligne est le bon indicateur à réguler : une ligne \emph{trop
large} sature le capteur et fausse le calcul du centre (centroïde)~; une ligne
\emph{trop fine} disparaît par endroits et crée des trous dans le nuage.
Réguler la largeur revient donc à maximiser la qualité du signal d'entrée de
tout le reste de la chaîne. La boucle garde le laser éteint par défaut et le
coupe en cas d'erreur, conformément à la politique de sécurité.

\subsection{Extraction de la ligne laser}

L'extraction transforme une image couleur en une liste de points sous-pixel
décrivant la ligne laser. C'est l'étape la plus en amont, donc la plus
critique pour le bruit final.

\paragraph{Algorithme.}
\begin{enumerate}
  \item \textbf{masque de région d'intérêt par caméra} : un masque polygonal
        écarte les bords du champ (supports, parois, plateau) où la ligne ne
        peut pas physiquement apparaître~;
  \item \textbf{seuillage sur le contraste vert} $G-R$ : le laser étant vert
        (520\,nm), la différence entre canaux vert et rouge isole la ligne tout
        en rejetant les reflets blancs et la lumière ambiante~;
  \item \textbf{centroïde sous-pixel par colonne} : pour chaque colonne de
        l'image, le centre lumineux est estimé en pondérant les intensités, ce
        qui donne une précision inférieure au pixel~;
  \item \textbf{filtre de longueur minimale} (\texttt{min\_line\_pixels}) : les
        segments trop courts, typiquement du bruit, sont rejetés.
\end{enumerate}

\paragraph{Évolution et justification.}
Le prototype intermédiaire se limitait à un seuillage $G-R$ global. La version
finale ajoute le masque par caméra et le filtre de longueur, deux ajouts rendus
nécessaires par le passage à deux caméras : chaque caméra voit une scène
différente (supports, angles, parois), et un masque dédié évite que ces
éléments parasites soient pris pour la ligne. Le choix du contraste $G-R$ plutôt
qu'un seuil de luminosité est justifié par la longueur d'onde verte du laser,
sur laquelle les capteurs CMOS (matrice de Bayer, deux fois plus de pixels
verts) sont les plus sensibles. Ces seuils sont réglables et visualisables en
direct depuis le mode ingénierie (\autoref{fig:ui-manuel}).

\subsection{Filtrage des points aberrants}

Après triangulation, chaque profil contient quelques points isolés
(réflexions, bruit de centroïde). Un filtre par \textbf{$k$ plus proches
voisins} s'appuyant sur un \texttt{KDTree} (SciPy) supprime les points dont le
voisinage local est trop clairsemé pour appartenir à une surface réelle.

\paragraph{Raison d'être.}
Cette étape est volontairement légère (adaptée au Raspberry Pi) et placée
\emph{avant} la reconstruction : un point aberrant non filtré tire localement
la surface reconstruite et crée une pointe, particulièrement visible après
Poisson. Filtrer en amont est moins coûteux que corriger le maillage en aval.

\subsection{Fusion des profils}

La fusion combine les profils acquis en un nuage unique et cohérent. Deux
mécanismes complémentaires sont employés et détaillés au chapitre Software :

\begin{itemize}
  \item \textbf{fusion des deux caméras} dans un repère plateau commun :
        couverture maximale dans les zones vues par une seule caméra,
        moyennage conservateur dans les zones de recouvrement pour éviter une
        double surface~;
  \item \textbf{fusion demi-tour} : un profil à $\theta$ et son symétrique à
        $\theta+180^\circ$ observent partiellement la même matière~; les
        doublons sont moyennés localement, ce qui réduit le bruit sans perdre
        de couverture.
\end{itemize}

L'enjeu central est le \emph{repère commun} : sans extrinsèques correctes, deux
nuages individuellement justes se superposent avec un décalage. C'est pourquoi
le diagnostic compare les nuages des deux caméras \emph{avant} fusion.

\subsection{Reconstruction du maillage : de l'enveloppe convexe à Poisson}

C'est l'évolution algorithmique la plus structurante du projet. Le prototype
intermédiaire fermait le nuage par une \textbf{enveloppe convexe} (trimesh).
Simple et robuste, cette méthode a une limite de principe : elle ne peut
représenter \emph{aucune} concavité et arrondit toute forme creuse, ce qui
faussait systématiquement les pièces non convexes.

La version finale repose sur la reconstruction de \textbf{Poisson} (Open3D),
encadrée par plusieurs étapes de préparation et de nettoyage :

\begin{enumerate}
  \item \textbf{estimation des normales} orientées de manière cohérente, sans
        laquelle Poisson ne peut pas reconstruire de surface~;
  \item \textbf{reconstruction de Poisson} (profondeur d'octree configurable,
        8 par défaut), qui produit une surface fermée et lisse à partir du
        nuage orienté~;
  \item \textbf{suppression des sommets de faible densité}, qui correspondent
        aux zones extrapolées par Poisson hors du nuage réel~;
  \item \textbf{détection et découpe du plan supérieur} : une bande planaire
        dense est détectée au sommet et le maillage est coupé par ce plan, ce
        qui supprime l'artefact en « dôme » typique de Poisson sur un nuage
        ouvert vers le haut~;
  \item \textbf{fermeture du fond} par un cap plat horizontal, l'axe vertical
        du plateau servant de référence.
\end{enumerate}

\paragraph{Raison d'être de chaque étape.}
Poisson est choisi parce qu'il gère les concavités et rend un maillage
\emph{étanche}, exploitable pour l'impression 3D. Mais appliqué seul sur un
nuage ouvert (l'objet n'est pas scanné par le dessus ni par dessous), il
« referme » la surface en inventant un dôme et un fond. Les étapes 3 à 5
existent précisément pour retirer ces inventions : on supprime ce qui est peu
dense, on coupe le dôme par le plan supérieur réellement mesuré, puis on
referme proprement le fond par une surface plane connue. C'est cette
combinaison qui a corrigé la déformation en hauteur observée à mi-parcours.

\subsection{Synchronisation moteur / laser / caméra}

La séquence d'acquisition est orchestrée par la machine d'états. Pour chaque
profil, l'ordre est~:

\begin{enumerate}
  \item rotation du plateau d'un pas~;
  \item stabilisation mécanique~;
  \item allumage du laser~;
  \item capture \textbf{des deux caméras}~;
  \item extinction du laser.
\end{enumerate}

La capture séquentielle des deux caméras au sein du même pas garantit qu'elles
observent rigoureusement la même position du plateau, condition nécessaire à
une fusion cohérente.

\subsection{Stabilité, précision et sources d'erreur}

Les contrôles manuels sur les scans réels indiquent des dimensions cohérentes
en largeur et en profondeur. La déformation en hauteur observée au rapport
intermédiaire a été \textbf{corrigée} par la découpe du plan supérieur et la
fermeture du fond. La précision absolue chiffrée reste à valider formellement.

\begin{table}[H]
  \centering
  \begin{tabular}{@{}p{4.8cm}p{4cm}p{4.5cm}@{}}
    \toprule
    \textbf{Source d'erreur} & \textbf{Réponse apportée} & \textbf{Statut} \\
    \midrule
    Occultation de la ligne laser & Seconde caméra + fusion calibrée & Traité \\
    Saturation / perte de la ligne & Régulation d'exposition automatique & Traité \\
    Points aberrants & Filtrage $k$-NN (KDTree) & Traité \\
    Déformation en hauteur & Découpe plan supérieur + cap de fond & Corrigé \\
    Calibration caméra/extrinsèques & Protocole de calibration par caméra & À fiabiliser \\
    Glissement de la pièce & Surface adhérente + centrage (méca) & À implémenter \\
    \bottomrule
  \end{tabular}
  \caption{Sources d'erreur et réponses algorithmiques ou mécaniques apportées.}
  \label{tab:erreurs}
\end{table}

\subsection{Validation prévue}

La validation métrologique devra s'appuyer sur des objets étalons imprimés en
3D ou mesurés précisément. Le protocole prévu consiste à scanner plusieurs fois
le même objet, à mesurer les dimensions reconstruites et à comparer la
répétabilité entre scans, afin de vérifier l'exigence \textbf{E24} sans se
limiter à une appréciation visuelle du STL.
